\begin{thebibliography}{99}
    \item A.V.エイホ, M.S.ラム, R.セシィ, and J.D.ウルマン 『コンパイラ: 原理・技法・ツール』 サイエンス社, 2009.第2版.
    \item 新屋良磨, 鈴木勇介, and 高田謙 『正規表現技術入門: 最新エンジン実装と理論的背景』 技術評論社, 2015.
    \item J.ホップクロフト, R.モトワニ, and J.ウルマン 『オートマトン 言語理論 計算論 I』 サイエンス社, 2003. 第2版.
    \item J.ホップクロフト, R.モトワニ, and J.ウルマン 『オートマトン 言語理論 計算論 II』 サイエンス社, 2003. 第2版.
    \item 中田育男 『コンパイラの構成と最適化』 朝倉書店, 2009. 第2版.
    \item Andrew W.Appel. \textit{Modern Compiler Implementation in ML}. Cambridge University Press, 1998.
    \item Terence Parr. \textit{Language Implementation Patterns}. Pragmatic Bookshelf, 2009.
    \item Terence Parr. \textit{The Definitive ANTLR 4 Reference}. Pragmatic Bookshelf, 2nd edition, 2013.
    \item Dick Grune and Ceriel J.H.Jacobs. \textit{Parsing Techniques: A Practical Guide}. Springer, 2nd edition, 2008.
    \item Keith D.Cooper and Linda Torczon. \textit{Engineering a Compiler}. Morgan Kaufmann, 2nd edition, 2011.
    \item Noam Chomsky. Three models for the description of language. \textit{IRE Transactions on Information Theory}, 2(3):113–124, 1956.
    \item Philip M.Lewis II and Richard E.Stearns. Syntax-directed transduction. \textit{Journal of the ACM}, 15(3):465–488, 1968.
    \item Donald E.Knuth. On the Translation of Languages from Left to Right. \textit{Information and Control}, 8(6):607–639, 1965.
    \item Jay Earley. An efficient context-free parsing algorithm. \textit{Communications of the ACM}, 13(2):94–102, 1970.
    \item Frank DeRemer. Simple LR(k) grammars. \textit{Communications of the ACM}, 14(7):453–460, 1971.
    \item Frank DeRemer and Thomas J. Pennello. Efficient Computation of LALR(1) Look-Ahead Sets. \textit{ACM Transactions on Programming Languages and Systems}, 4(4):615–649, 1982.
    \item Masaru Tomita. \textit{Efficient Parsing for Natural Language}. Kluwer Academic Publishers, 1985. GLR 法の原典.
\clearpage
    \item Bryan Ford. Packrat parsing: simple, powerful, lazy, linear time. In \textit{ICFP ’02}, pages 36–47, 2002.
    \item Bryan Ford. Parsing expression grammars: a recognition-based syntactic foundation. In \textit{POPL ’04}, pages 111–122, 2004.
    \item Alessandro Warth, James R. Douglass, and Todd Millstein. Packrat Parsers Can Support Left Recursion. In \textit{PEPM ’08}, pages 103–110, 2008.
    \item Elizabeth Scott and Adrian Johnstone. GLL Parsing. \textit{Electronic Notes in Theoretical Computer Science}, 253(7):177–189, 2010.
    \item Terence Parr and Kathleen Fisher. LL(*): The Foundation of the ANTLR Parser Generator. In \textit{PLDI ’11}, pages 425–436, 2011.
    \item Terence Parr, Sam Harwell, and Kathleen Fisher. Adaptive LL(*) Parsing: The Power of Dynamic Analysis. In \textit{OOPSLA ’14}, pages 579–598, 2014.
    \item Tim A.Wagner and Susan L.Graham. Efficient and Flexible Incremental Parsing. \textit{ACM Transactions on Programming Languages and Systems}, 20(5):980–1013, 1998.
    \item Lukas Diekmann and Laurence Tratt. Don’t Panic! Better, Fewer, Syntax Errors for LR Parsers. In \textit{ECOOP ’20}, pages 6:1–6:32, 2020.
    \item ECMA International. ECMA-404 The JSON Data Interchange Syntax, 2017. \url{https://ecma-international.org/publications-and-standards/standards/ecma-404/}.
    \item Tim Bray. RFC 8259: The JavaScript Object Notation (JSON) Data Interchange Format, 2017. \url{https://datatracker.ietf.org/doc/html/rfc8259}.
    \item ISO/IEC. ISO/IEC 14977:1996 Extended BNF, 1996. \url{https://www.iso.org/standard/26153.html}.
    \item Dave Crocker and Paul Overell. RFC 5234: Augmented BNF for Syntax Specifications: ABNF, 2008. \url{https://datatracker.ietf.org/doc/html/rfc5234}.
    \item Pete Resnick. RFC 5322: Internet Message Format, 2008. \url{https://datatracker.ietf.org/doc/html/rfc5322}.
    \item Pablo Galindo, Lysandros Nikolaou, and Guido van Rossum. PEP 617: New PEG parser for CPython, 2020. \url{https://peps.python.org/pep-0617/}.
    \item Robert Nystrom. Crafting Interpreters. \url{https://craftinginterpreters.com/}.
    \item Jack Crenshaw. Let’s Build a Compiler. \url{https://compilers.iecc.com/crenshaw/}.
    \item Thorsten Ball. Writing An Interpreter In Go. \url{https://interpreterbook.com/}.
    \item ANTLR. ANTLR Official Site. \url{https://www.antlr.org/}.
    \item JavaCC. JavaCC Official Site. \url{https://javacc.github.io/javacc/}.
    \item GNU Project. GNU Bison Manual. \url{https://www.gnu.org/software/bison/manual/}.
    \item tree-sitter. Tree-sitter Official Site. \url{https://tree-sitter.github.io/tree-sitter/}.
\end{thebibliography}